\documentclass{article}
\usepackage{mathtools} 
\usepackage{fontspec}
\usepackage[UTF8]{ctex}
\usepackage{amsthm}
\usepackage{mdframed}
\usepackage{xcolor}
\usepackage{amssymb}
\usepackage{amsmath}


% 定义新的带灰色背景的说明环境 zremark
\newmdtheoremenv[
  backgroundcolor=gray!10,
  % 边框与背景一致，边框线会消失
  linecolor=gray!10
]{zremark}{说明}

% 通用矩阵命令: \flexmatrix{矩阵名}{元素符号}{行数}{列数}
\newcommand{\flexmatrix}[4]{
  \[
  #1 = \begin{pmatrix}
    #2_{11}     & #2_{12}     & \cdots & #2_{1#4}   \\
    #2_{21}     & #2_{22}     & \cdots & #2_{2#4}   \\
    \vdots      & \vdots      & \ddots & \vdots     \\
    #2_{#31}    & #2_{#32}    & \cdots & #2_{#3#4}
  \end{pmatrix}
  \]
}

% 简化版命令（默认矩阵名为A，元素符号为a）: \quickmatrix{行数}{列数}
\newcommand{\quickmatrix}[2]{\flexmatrix{A}{a}{#1}{#2}}

\begin{document}
\title{习题 4.2}
\author{张志聪}
\maketitle

\section*{6}

 (1) $f$在$[a, +\infty)$上有界。

因为$\lim\limits_{x \to +\infty} f(x)$存在，
由局部有界性可知，$\exists G > 0, M > 0$，当$x > M$时，有
\begin{align*}
  |f(x)| \leq G
\end{align*}
$\because f$在$[a, +\infty)$上连续 \\
$\therefore f$在$[a, M]$上连续，由有界性定理可知，\\
$\exists P > 0$，对$\forall x \in [a, M]$，有
\begin{align*}
  |f(x)| \leq P
\end{align*}
取$Q = max(P, G)$，对$\forall x \in [a, +\infty)$有
\begin{align*}
  |f(x)| \leq Q
\end{align*}
综上， $f$在$[a, +\infty)$上有界。

(2) $f$在$[a, +\infty)$上必有最大值或最小值。

设$\lim\limits_{x \to +\infty} f(x) = L$。

(2.1)如果$\exists x_0 \in [a, +\infty)$，有$f(x_0) > L $，
由保号性可知，$\exists M > x_0$，当$x > M$时，有
\begin{align*}
  f(x) < f(x_0)
\end{align*}
$\because f$在$[a, +\infty)$上连续 \\
$\therefore f$在$[a, M]$上连续，由最大、最小值定理，
$f$在$[a, M]$上有最大值，不妨设为$P$。\\
$\because x > M$时，有
\begin{align*}
  f(x) < f(x_0) < P
\end{align*}
所以，$P$就是$f$在$[a, +\infty)$上的最大值。

(2.2)如果$\exists x_0 \in [a, +\infty)$，有$f(x_0) < L $，
由保号性可知，$\exists M > x_0$，当$x > M$时，有
\begin{align*}
  f(x_0) < f(x)
\end{align*}
$\because f$在$[a, +\infty)$上连续 \\
$\therefore f$在$[a, M]$上连续，由最大、最小值定理，
$f$在$[a, M]$上有最小值，不妨设为$p$。\\
$\because x > M$时，有
\begin{align*}
  p < f(x_0) < f(x)
\end{align*}
所以，$p$就是$f$在$[a, +\infty)$上的最小值。

(2.3)如果$!\exists x_0 \in [a, +\infty)$，有$f(x_0) < L$或$f(x_0) > L$，
那么，$f(x) = L$，即$f(x)$是常值函数，
此时，$L$既是$f$在$[a, +\infty)$上的最大值，也是$f$在$[a, +\infty)$上的最小值。

综上，$f$在$[a, +\infty)$上必有最大值或最小值。

\section*{13}

\begin{itemize}
  \item (1) 在$[a, b]$上一致连续。

        $\because f$在$[a, b]$上连续，\\
        $\therefore$由一致连续性定理，$f$在$[a, b]$上一致连续。

  \item (2) 在$(-\infty, +\infty)$上不一致连续。

        可取$\epsilon = 1$，对无论多么小的正数$\delta$，取$x' = n + \frac{\delta}{2}, x'' = n$，
        总有
        \begin{align*}
          \lim\limits_{n \to +\infty} |f(x') - f(x'')|
           & = \lim\limits_{n \to +\infty} |(n + \frac{\delta}{2})^2 - n^2| \\
           & = \lim\limits_{n \to +\infty} |n\delta + \frac{\delta^2}{4}|   \\
           & = + \infty
        \end{align*}
        于是，由$\lim\limits_{n \to +\infty} |f(x') - f(x'')| = +\infty$，
        $\exists N > 0 $，当$x > N$时，有
        \begin{align*}
          |f(x') - f(x'')| > 1
        \end{align*}
        所以，$f$在$(-\infty, +\infty)$上不一致连续。
\end{itemize}

\section*{16}

$\because \lim\limits_{x \to +\infty} f(x)$存在，由柯西准则可知，
$\forall \epsilon > 0, \exists M > a$，当$x', x'' > M$时，有
\begin{align*}
  |f(x') - f(x'')| < \epsilon
\end{align*}
$\therefore$对任意$\delta > 0$，存在$x', x'' \in [M + 1, +\infty)$且$|x' - x''| < \delta$，有
\begin{align*}
  |f(x') - f(x'')| < \epsilon
\end{align*}
$\therefore$ $f$在$[M + 1, +\infty)$上一致连续。

$\because f$在$[a, +\infty)$上连续，
$\therefore$由一致连续性定理，$f$在$[a, M + 1]$上一致连续。

综上，$f$在$[a, +\infty)$上一致连续。

\section*{17}

提示
\begin{align*}
  g(x) = f(x + a) - f(x)
\end{align*}
考虑$g(x) = 0$的解。

\end{document}